\cleardoublepage
\phantomsection
\addcontentsline{toc}{chapter}{Sommario}
\chapter*{Sommario}

L'annotazione manuale di riquadri di delimitazione (\textit{bounding box}) per l'object detection richiede tempo e risorse considerevoli, in particolare nei domini specialistici come i beni culturali. In questo lavoro si esamina il paradigma della \textit{Weakly-Supervised Object Localization} (WSOL) applicato al dataset egocentrico UNICT-VEDI POI, che raccoglie 57 classi di elementi monumentali nel Monastero dei Benedettini di Catania.

Sfruttando le feature di auto-segmentazione del modello DINOv3 ViT-S/16 con testa lineare (\textit{linear probing}), è stata definita una procedura basata su Linear-CAM per ricavare pseudo-annotazioni spaziali a partire dalle sole etichette di classe dell'immagine. Le etichette così generate servono ad addestrare il rilevatore compatto YOLOv8n, senza impiegare coordinate cartesiane supervisionate nel processo di apprendimento.

Per quantificare il risparmio di annotazione, è stato impostato un confronto di \textit{sample efficiency} rispetto a baseline supervisionate addestrate su quote progressive ($1\% \dots 100\%$) del ground truth con tre repliche stocastiche. Il modello distillato ottiene il $37.39\%$ di mAP@50, prestazione equivalente all'$1.43\%$ di dati etichettati a mano e pari al $61.79\%$ della capacità del rilevatore supervisionato al $100\%$. L'analisi per scala dimensionale evidenzia una ritenzione del $63.60\%$ nel Tier Macro (26 classi di taglia intermedia non degeneri) e dell'$84.14\%$ nel Tier Full-Frame. L'esame dei fallimenti nel Tier Tiny ($0.00\%$ mAP@50) mostra inoltre che l'errore sistematico non deriva dalla risoluzione spaziale delle patch, ma dal bias di contesto e dall'\textit{attribution leakage}: i massimi di attivazione del classificatore scivolano sulle strutture murarie contigue anziché sul piccolo reperto da localizzare.
